\section*{第 \ref{ch:b3:plant-molecular-physiology} 章——植物分子生理学与胁迫响应}

\begin{solution}{exo:b3:plant-molecular-physiology:1}
\omterm{def:b3:plant-molecular-physiology:photoreceptors}{光敏色素}，红光 \qty{660}{nm} 与远红光 \qty{730}{nm}：
去\omterm{def:b3:plant-molecular-physiology:photoreceptors}{黄化}、\omterm{prop:b3:plant-molecular-physiology:photoequilibrium}{避荫}、萌发。\omterm{def:b3:plant-molecular-physiology:photoreceptors}{隐花色素}，蓝光 \qty{450}{nm}：
去\omterm{def:b3:plant-molecular-physiology:photoreceptors}{黄化}、生物钟、开花。\omterm{def:b3:plant-molecular-physiology:photoreceptors}{向光素}，蓝光：
向光弯曲、\omterm{prop:b3:plant-molecular-physiology:stomata}{气孔}开放、叶绿体运动。
叶绿素把光当作能量来测量，却
不利用其中的信息；光感受器读取颜色、方向和时长，
所用通量低一百万倍左右，并借此改变基因表达——
这正是“吃掉光”与“读取光”的区别。
\end{solution}

\begin{solution}{exo:b3:plant-molecular-physiology:2}
每种\omterm{def:b3:endocrinology:hormone}{激素}都促使一种阻遏蛋白泛素化并被降解：
\omterm{def:b3:plant-molecular-physiology:hormones}{生长素}将 \omterm{def:b3:plant-molecular-physiology:hormones}{TIR1} 与 Aux/IAA 蛋白结合在一起，\omterm{def:b3:plant-molecular-physiology:hormones}{赤霉素}与 GID1 结合后
将 \omterm{def:b3:plant-molecular-physiology:hormones}{DELLA} 蛋白递送至 F-box 连接酶，\omterm{def:b3:plant-molecular-physiology:immunity}{茉莉酸}将 COI1 与 JAZ 蛋白结合在一起。
阻遏蛋白原本就结合在靶基因上，激活蛋白则已在旁待命，
因此无需合成新的蛋白质：
蛋白酶体在数分钟内降解阻遏蛋白，
靶基因便立即开启。
\end{solution}

\begin{solution}{exo:b3:plant-molecular-physiology:3}
\omterm{def:b3:plant-molecular-physiology:hormones}{脱落酸}与 PYR/PYL 结合；该复合体抑制 PP2C 磷酸酶；
SnRK2 激酶（OST1）不再被去磷酸化，随即活化，并
使 SLAC1 阴离子通道磷酸化；阴离子外流，膜
去极化，外向钾离子通道开放，K$^{+}$ 随后
与水分一同外流，膨压下降，\omterm{prop:b3:plant-molecular-physiology:stomata}{保卫细胞}一同萎蔫。任何一个环节均可
被破坏：缺少受体，磷酸酶无法被抑制（经典的
\emph{abi1} 突变体）、缺少 OST1 或缺少 SLAC1——每种情况都会使
植株萎蔫，且干旱时\omterm{prop:b3:plant-molecular-physiology:stomata}{气孔}仍保持开放。
\end{solution}

\begin{solution}{exo:b3:plant-molecular-physiology:4}
PTI 识别整类微生物所共有的分子
（如鞭毛蛋白、几丁质），依赖细胞表面的受体激酶，
其结局是细胞壁加固、\omterm{prop:b3:plant-molecular-physiology:stomata}{气孔}关闭、防御基因表达，生长减缓，
通常不伴随细胞死亡。ETI 识别特定的
效应蛋白或其造成的\omterm{def:b3:dna-repair:lesions}{损伤}，借助胞内 NLR 受体进行识别，
其结局是被感染细胞发生超敏性死亡，并启动系统
性警报。
\end{solution}

\begin{solution}{exo:b3:plant-molecular-physiology:5}
日光 $0.87\times 1.15/1.75 = 0.57$；单片叶片 $0.87\times 0.4/1.0 =
0.35$；密集冠层 $0.87\times 0.1/0.7 = 0.12$；白炽灯
$0.87\times 0.7/1.3 = 0.47$。白炽灯富含远红光，
因此 $\varphi$ 低于日光，幼苗感知到的是部分遮荫：
幼苗随之伸长。冷白灯与不含远红光的 LED 则
产生相反的效果。
\end{solution}

\begin{solution}{exo:b3:plant-molecular-physiology:6}
细胞壁 $5.0$，细胞溶胶 $7.5$：$(1 + 10^{2.75})/(1 + 10^{0.25}) = 563/2.78
= 203$。细胞壁 $4.5$：$563/(1 + 10^{-0.25}) = 563/1.56 = 360$——
细胞壁酸化提高了外侧中性形式的比例，也提高了吸收量。
细胞溶胶 $6.5$、细胞壁 $5.5$：$(1 + 10^{1.75})/6.62 = 57/6.62 = 8.6$：
滞留量降至约五分之一，\omterm{def:b3:plant-molecular-physiology:hormones}{生长素}从各个面渗漏出去，极性
梯度随之变平——这也是缺氧根系失去方向性的原因之一。
\end{solution}

\begin{solution}{exo:b3:plant-molecular-physiology:7}
\qty{1}{cm/h} $= \qty{10000}{\micro m/h}$：每小时 $200$ 个细胞，
每个细胞 \qty{18}{s}。穿过 \qty{50}{\micro m} 的扩散：$t =
(5\times 10^{-5})^{2}/(2\times 5\times 10^{-10}) = \qty{2.5}{s}$。
细胞内部只需数秒即可穿过；大部分 \qty{18}{s} 都耗在
等待经基部膜上的 PIN 蛋白将其运出，
这才是限速步骤。
\end{solution}

\begin{solution}{exo:b3:plant-molecular-physiology:8}
$E/A = 1.6\,\Delta w/\Delta c = 1.6\times\num{16000}/150 = 171$ 个水分子，
每个 CO$_{2}$。C$_{4}$ 植物，$\Delta c = 250$：$102$。干燥的白天，
$\Delta w = 24$：$256$。
\end{solution}

\begin{solution}{exo:b3:plant-molecular-physiology:9}
冷冻通过将水分拉入细胞外的冰中而使细胞脱水，因此
低温与干旱都属于脱水胁迫。二者共享第二信使
（ABA、钙信号峰、活性氧）、转录因子
（CBF/DREB 家族可被二者同时诱导），以及产物：
脱水素和 LEA 蛋白，它们保护膜与蛋白质免于失水，
相容性溶质（脯氨酸、糖类）可保持水分，
此外还有抗氧化剂。一株植物若已为某一种胁迫产生了这些物质，
便也具备应对另一种胁迫的能力。
\end{solution}

\begin{solution}{exo:b3:plant-molecular-physiology:10}
在全日照的 \qty{1}{\%} 时，$k_{1} + k_{2} = \qty{0.005}{s^{-1}}$：达到平衡的 95~\%
所需时间为 $3/k = \qty{600}{s}$，即十分钟，而正午仅需六
秒。八小时相当于四个半衰期：$1/16 \approx
\qty{6}{\%}$ 的 Pfr 残留。短暂闪光只能设定 $\varphi$，但
随后 Pfr 便衰减，而整天的光照使其在数小时内保持较高水平；
需要 Pfr 持续作用数小时的响应（萌发、去\omterm{def:b3:plant-molecular-physiology:photoreceptors}{黄化}）会对
曝光量进行积分，因此被犁一时翻起的种子
不会像躺在地表的种子那样作出反应。
\end{solution}

\begin{solution}{exo:b3:plant-molecular-physiology:11}
每天十个位点，每个位点 $100$ 个细胞：共 \num{1000} 个细胞，占叶片的 $10^{-4}$
；整个季节下来仅占不足百分之一，而
若一次感染不受控制地扩散，则会损失整片叶片。死亡细胞几乎
没有代价，还会带走病原体的食物——对于活体营养型病原体，
它需要活细胞。\omterm{def:b3:cell-cycle-apoptosis:apoptosis}{坏死}营养型病原体以死亡组织为食：\omterm{def:b3:plant-molecular-physiology:immunity}{超敏反应}
为其提供了食物与立足点，而其中一些
（\emph{Botrytis}、\emph{Sclerotinia}）会分泌毒素，并有意诱发该反应；
对付它们时，植株依赖的是\omterm{def:b3:plant-molecular-physiology:immunity}{茉莉酸}介导的防御，
而非细胞死亡。
\end{solution}

\begin{solution}{exo:b3:plant-molecular-physiology:12}
以速率 $p$ 从第 12 小时持续到第 16 小时生产。长日照，全程有光，
$k = \ln 2/4 = \qty{0.17}{h^{-1}}$：第 16 小时的水平为 $(p/k)(1 -
\mathrm{e}^{-4k}) = (p/0.17)(0.5) = 2.9p$。短日照，自第
10 小时起转入黑暗，$k = \qty{1.39}{h^{-1}}$：$(p/1.39)(1 - \mathrm{e}^{-5.5})
\approx 0.72p$。长日照下 CONSTANS 多出四倍。一个阈值
介于——比如 $2p$，即激活 \emph{FT} 所需的量——之间，
只有当时钟的生成时段与光照重合时才被越过：
内部节律与外部昼夜的重合，把渐变的日长
转化为全或无式的开花决定。
\end{solution}

\begin{solution}{pb:b3:plant-molecular-physiology:1}
\textbf{1.} 日光下 $0.87\times 1.15/1.75 = 0.57$；冠层下 $0.87\times
0.2/0.8 = 0.22$。
\textbf{2.} $r = 2(1 - \varphi)$：日光下为 \qty{0.86}{mm/h}，
荫蔽下为 \qty{1.57}{mm/h}；\qty{48}{h} 内为 $41$，对比 \qty{75}{mm}，
多出 \qty{34}{mm}。
\textbf{3.} 红光降至远红光的 $0.1\times 1.15 = 0.115$，远红光
降至 $0.8$：$\zeta = 0.144$，$\varphi = 0.87\times 0.144/0.744 =
0.17$。是，单片叶使 $\varphi$ 由 $0.57$ 降至 $0.17$，
处于避阴反应范围的深处。
\textbf{4.} 两种光转换速率都与通量成正比，因此它们的
比值不变。阴天时光谱几乎不变，
$\varphi$ 保持在 $0.57$ 附近：阴天时空旷处的幼苗不会
伸长——它感知的是邻株，而非亮度。
\textbf{5.} \qty{95}{\%} 所需时间为 $3/(k_{1} + k_{2})$：全日照下 \qty{6}{s}，
在 \qty{1}{\%} 光照下为 \qty{600}{s}。
\textbf{6.} \qty{8}{h} $=$ 四个半衰期，$1/16 = \qty{6.3}{\%}$；
\qty{14}{h} $=$ 七个，$1/128 = \qty{0.8}{\%}$。黎明时残留的 Pfr
可以反映夜长，但逆转是热反应，
在温暖的夜晚进行得更快，而起始值取决于
白天的光照，植物转而借助时钟来测量夜长。
\textbf{7.} 细胞壁：$1/(1 + 10^{0.75}) = 0.15$；细胞溶胶：$1/(1 +
10^{2.45}) = 0.0035$。
\textbf{8.} $(1 + 282)/(1 + 5.6) = 283/6.6 = 43$。细胞溶胶
中 $43\times 0.1 = \qty{4.3}{\micro M}$。
\textbf{9.} $\qty{1}{cm/h}/\qty{100}{\micro m} = 100$ 个细胞每小时，
每个耗时 \qty{36}{s}。
\textbf{10.} \qty{5}{h}。一小时内即开始弯曲，因为
真正起作用的再分配是侧向的，横跨尖端一毫米，
反应是局部的；长距离运输流决定的是供应，
而非时机。
\textbf{11.} 下侧：$60/50 = 1.2$，\omterm{def:b3:plant-molecular-physiology:hormones}{生长素}高出常规 \qty{20}{\%}，
伸长率 $0.8$；上侧\omterm{def:b3:plant-molecular-physiology:hormones}{生长素}为常规的 $0.8$，伸长率 $1.2$。上
侧生长快于下侧：根向下弯曲。
\textbf{12.} 茎的细胞位于其\omterm{def:b3:plant-molecular-physiology:hormones}{生长素}最适浓度之下，因此额外的
\omterm{def:b3:plant-molecular-physiology:hormones}{生长素}在下侧促进其生长：下侧
生长快于上侧，茎向上弯曲。
\textbf{13.} $E = 0.3\times 16 = \qty{4.8}{mmol\,m^{-2}\,s^{-1}}$。
\textbf{14.} $A = (0.3/1.6)\times 150 = \qty{28}{\micro mol\,m^{-2}\,s^{-1}}$；
$E/A = 4800/28 = 171$。
\textbf{15.} 面积 \qty{0.002}{m^2}，时间 \num{43200}~s：水分 $4.8\times
10^{-3}\times 0.002\times\num{43200} = \qty{0.41}{mol} = \qty{7.5}{g}$；
CO$_{2}$ $28\times 10^{-6}\times 0.002\times\num{43200} =
\qty{2.4e-3}{mol}$，即 $2.4\times 10^{-3}/6\times 180 =
\qty{0.073}{g}$ 葡萄糖——即耗水一百克，
换得一克糖。
\textbf{16.} $E = 0.05\times 16 = \qty{0.8}{mmol\,m^{-2}\,s^{-1}}$，
$A = \qty{4.7}{\micro mol\,m^{-2}\,s^{-1}}$，比值仍为 $171$。植株
节省了其水分的七分之六，并放弃了其
碳的七分之六：关闭阀门改变的是速率，而非代价。
\textbf{17.} C$_{4}$：$\Delta c = 250$，$E/A = 4800/46.9 = 102$。CAM 在
夜间：$E = 0.3\times 5 = \qty{1.5}{mmol}$，$A = 28$：$E/A = 53$。
\textbf{18.} $E/A = 1.6\,\Delta w/\Delta c$：电导约去，
因为两种气体通过的是同一\omterm{prop:b3:plant-molecular-physiology:stomata}{气孔}。植物可以改变
梯度——在内部富集 CO$_{2}$（C$_{4}$），在空气
湿润时开放（CAM，黎明），降低叶温，用绒毛增厚边界层
——但改变不了两种气体共享一个孔的物理规律。
\textbf{19.} 每天 \num{1000} 个细胞，占叶片的 $10^{-4}$。
\textbf{20.} 60 天内 \num{60000} 个细胞，占 \qty{0.6}{\%}。若每
天逃逸一次并摧毁 \qty{5}{\%}，二十天内即可耗尽整片叶片：
\omterm{def:b3:plant-molecular-physiology:immunity}{超敏反应}的代价只是单次逃逸感染
代价的百分之一。
\textbf{21.} $0.2\times 5 = \qty{1}{\%}$ 的光合作用。若始终
维持，则将耗费 \qty{5}{\%}，植株将被
那些将其投入叶片的邻株超过；诱导防御只有在
威胁存在时才划算。
\textbf{22.} 它获得了对该 NLR 的隐蔽性，并感染
抗性品种；它失去了效应蛋白原有的作用（抑制 PTI），
因此在不含该 NLR 的植株上较弱。植物种群逐渐偏向
针对其余效应蛋白的 NLR，旧的抗性
基因失去作用而减少：双方
基因的频率依赖性循环。
\textbf{23.} 拟态物（冠菌素）激活\omterm{def:b3:plant-molecular-physiology:immunity}{茉莉酸}途径，
该途径拮抗\omterm{def:b3:plant-molecular-physiology:immunity}{水杨酸}途径；\omterm{def:b3:plant-molecular-physiology:immunity}{水杨酸}是对
活体营养型病原体（如细菌本身）的防御，而拟态物还
重新打开了植株已关闭的\omterm{prop:b3:plant-molecular-physiology:stomata}{气孔}。病原体使免疫系统的一支
与另一支相互对抗。
\textbf{24.} 在数百万株相同植株上使用单一抗性基因，
就是一种均一选择：任何丢失或改变效应蛋白的突变体
会在几个季节内蔓延至整个作物群体。育种者会叠加多个
抗性基因，使多个效应蛋白须同时丢失，混合种植
不同品种、逐年轮换基因，并偏重
广谱的模式触发抗性，单一突变无法逃逸。
\textbf{25.} 冠层下 $\varphi \approx 0.22$；\omterm{def:b3:plant-molecular-physiology:hormones}{生长素}
内侧与外侧之比 $\approx 43$；约 $170$ 个水分子
随每固定一个 CO$_{2}$ 而散失。
\end{solution}
